\section{Toward Auditable Evaluation: A Minimal Provenance Standard}
\label{sec:provenance}

The clearest actionable output of this study is a recommendation to the community:
a benchmark harness should record the \emph{execution substrate}, not just the
recipe and the model name. Our investigation of HELM's own deployment metadata
quantifies the gap: of 148 \texttt{HuggingFaceClient} deployments in HELM's
\texttt{model\_deployments.yaml}, only 19 pin \texttt{torch\_dtype} (all
\texttt{bfloat16}) and only 7 pin a model \texttt{revision}. The parameters that most
change a greedy output are precisely the ones least often recorded
(Appendix~\ref{app:unrecorded}).

\paragraph{Three fields close most of the deployment-side gap.} The high-leverage
observation is that many parameters collapse into one:
\begin{quote}
\texttt{model\_revision} $+$ \texttt{torch\_dtype} $+$ \texttt{transformers\_version}
\end{quote}
The revision fixes the tokenizer, chat template, and generation config \emph{as of
then} (they all live inside the model repository); the \texttt{transformers} version
fixes both the precision default and whether the chat template honors
\texttt{add\_generation\_prompt}; the dtype pins the load precision directly. Each of
the OLMo findings (\S\ref{sec:olmo}) reduces to one of these three fields --- (a)
to the revision and (b) to the \texttt{transformers} version, both confirmed, and
(c) to the dtype, which remains a discovery.

We scope that claim deliberately. Of the parameters in
Table~\ref{tab:identifiability} that are not already identifiable, these three
address the deployment coordinate --- dtype, tokenizer behavior, chat
template, revision --- and leave the instrument coordinate untouched: engine
build and kernel versions, batching, cache, device topology, hardware, and
harness era.
Those are precisely the parameters that make a hosted endpoint nondeterministic
even at temperature zero (\S\ref{sec:related}), and no three data fields can pin
them; they require the container digest and engine record that the checklist below
asks for separately. The evidence for the three-field claim is likewise bounded: it
is the mapping above, on one model family, two of whose three legs are confirmed.
What we claim is that these three fields are the highest-leverage additions
available at negligible cost, not that recording them makes a run reproducible.

\paragraph{Record them out of the identity string.} These fields must live in a
machine-readable provenance block (HELM's existing recipe-facts slot suffices ---
no schema change), \emph{not} in the run-spec identity or the
\texttt{model\_deployment} name. A revision SHA in the pairing key would break
official$\leftrightarrow$local matching, whereas a SHA in the provenance block is
invisible to pairing and fully recoverable. Both weight-loading engines already
accept \texttt{revision} as a data-only field; the only code gap is a bundle
boundary that currently drops it.

\paragraph{A reproduction-target-aware release checklist.} Beyond the three
fields, a paper-grade evaluation artifact should preserve (Appendix~\ref{app:checklist}):
source commit and container digest; dataset and scorer versions; model and
tokenizer commit SHAs; chat-template text/hash and tokenized-prompt hash; dtype,
quantization, engine, attention implementation, and generation config; batching,
cache, tensor-parallel, and numerical/hardware settings; raw requests, token ids,
responses, metrics, and artifact-transformation history; \emph{every}
deployment-search candidate, not only the winner; and all analysis/figure scripts.
Crucially, the release should record, for each experiment, its reproduction
\emph{target} (\S\ref{sec:targets}) and three independent freshness facts ---
run-artifact, comparison, and report freshness --- so that a recent report
timestamp is never mistaken for an acceptance criterion (a lesson learned the hard
way when reports outlived the raw runs they summarized).
